\documentclass[11pt]{article}
\usepackage{amsmath,amssymb,amsthm}
\usepackage[margin=1.1in]{geometry}
\usepackage{booktabs,listings}
\usepackage[hidelinks]{hyperref}
\allowdisplaybreaks[2]
\urlstyle{same}

\newtheorem{theorem}{Theorem}[section]
\newtheorem{lemma}[theorem]{Lemma}
\newtheorem{corollary}[theorem]{Corollary}
\theoremstyle{definition}
\newtheorem{remark}[theorem]{Remark}
\newcommand{\Z}{\mathbb Z}
\newcommand{\Q}{\mathbb Q}

\title{Odd-Step Budgets for Collatz Excursions and\\
An Exact Certificate for 108 Local Minima}
\hypersetup{pdftitle={Odd-Step Budgets for Collatz Excursions and
An Exact Certificate for 108 Local Minima}}
% Supply the confirmed author names, affiliations, and email before submission.
\author{}
\date{}

\begin{document}
\maketitle

\begin{abstract}
Let \(c\) be a fixed positive odd integer, and iterate the map
\(n\mapsto n/2\) for even \(n\) and \(n\mapsto(3n+c)/2\) for odd \(n\).
For each fixed \(J\) and \(\eta>0\), the first \(J\) maximal odd runs
from an odd seed \(n\) contain at most
\((1+\eta)\log_2(n+c)+O_{c,J,\eta}(1)\) odd steps. The proof applies
the \(S\)-adic Subspace Theorem to a telescoping identity. A projection
on the same coordinates at all three places handles the exceptional
subspaces. It follows that the optimal power for the maximal height
of these excursions is \(\log_2 3\). For the classical map the budget
also yields \(M/\ln s\to\infty\) along hypothetical cycles with \(s\)
odd terms and \(M\) local minima. The same run identities give an
effective two-run height test using a binary window and its incoming
borrow. The test also exhibits a finite-scale instance of the
logarithmic excess in the height bound. Combined with unimodular
brackets, it proves that every nontrivial positive Collatz cycle
has at least \(108\) local minima, using the published verification
below \(2^{71}\) and the explicit logarithmic-form bound of Rhin.
The complete verifier uses integer and rational arithmetic.
\end{abstract}

\noindent\textit{2020 Mathematics Subject Classification.}
11J61, 11B83, 37P05, 11Y16.

\section{Introduction}

Fix a positive odd integer \(c\) and define the shortcut map on positive
integers by
\[
 S_c(x)=
 \begin{cases}
 x/2,& 2\mid x,\\
 (3x+c)/2,& 2\nmid x .
 \end{cases}
\]
Write \(S=S_1\) and \(\lambda=\log_2 3\). Starting at an odd seed \(n\),
an \emph{odd run} consists of consecutive odd inputs, up to the last
odd input before an even output. Between runs, at least one even step
occurs. Denote by \(A_{J,c}(n)\) the total number of odd steps in the
first \(J\) runs, and by \(H_{J,c}(n)\) the largest state from the seed
through the even output at the end of run \(J\). Each odd step increases
the state and each even step decreases it. Section~\ref{sec:budget}
shows that every run is finite, so these quantities are defined for
every \(J\ge 1\).

The elementary bound for \(H_{J,c}(n)+c\) has power \(\lambda^J\).
Theorem~\ref{thm:budget} replaces all \(J\) separate length bounds by
a single odd-step budget of slope \(1+\eta\). In particular,
Corollary~\ref{cor:height} identifies the optimal power \(\lambda\)
for every fixed \(J\). A family for \(c=1\) shows that for \(J\ge2\)
the pure power bound \(O(n^\lambda)\) fails.

For \(c=1\), Brox compressed a cycle into groups of rising steps
and proved finiteness with fewer than twice the logarithm of the
period many descents \cite{Brox}. Weaver later gave a fixed-descent
height bound for the odd-to-odd map and derived superlogarithmic
descent counts \cite{Weaver}. To compare the conventions, an odd run
of length \(a\) in the shortcut map comprises \(a-1\) odd-to-odd
steps with division exponent one, followed by one with exponent
at least two. Its even peak is at most twice the last odd input.
Thus Weaver's height estimate gives the same power for fixed
\(J\) when \(c=1\). Weaver also reports a stronger effective
\(\Omega(\ln s\ln\ln s)\) cycle bound \cite{Weaver}.
The proof here gives the odd-step budget directly from a forward
identity and applies to every positive odd \(c\).

To obtain an explicit cycle bound, we return to the same run
identities at two consecutive runs. Hercher proved \(M\ge92\)
\cite{Hercher}, and Wang's preprint gives \(M\ge96\) \cite{Wang}.
Alick's public Lean development combines Hercher's reduction with
two-rise bounds to prove \(M\ge102\) \cite{Alick}. We prove
\(M\ge108\) from the verification threshold \(2^{71}\) and Rhin's
bound \cite{Barina,SimonsWeger} by combining six unimodular brackets
with seventeen certified two-run height bounds. The logarithmic
excess in Remark~\ref{rem:excess} explains why the height caps
require more than a uniformly bounded additive margin in the
logarithmic scale. Section~\ref{sec:caps} gives a precise
finite-scale example and the congruence test that certifies the
chosen margins. Appendix~\ref{app:code} contains the complete code.

\section{An odd-step budget}\label{sec:budget}

\subsection{Run identities}

Let \(x_i\) be the start of run \(i\) and factor
\(x_i+c=2^{a_i}u_i\) with \(u_i\) odd. During this run,
\[
 S_c^j(x_i)+c=3^j2^{a_i-j}u_i \qquad(0\le j\le a_i).
\]
The state is odd for \(j<a_i\) and even for \(j=a_i\).
The even peak \(3^{a_i}u_i-c>0\) is followed by exactly \(v_i\ge1\)
even steps before \(x_{i+1}\). Consequently
\begin{equation}\label{eq:transition}
 2^{a_{i+1}+v_i}u_{i+1}=3^{a_i}u_i+c(2^{v_i}-1).
\end{equation}
Since \(a_i\le\log_2(x_i+c)\) and even steps decrease \(x+c\),
\begin{equation}\label{eq:elementary}
 \log_2(x_i+c)\le\lambda^{i-1}\log_2(n+c),\qquad
 A_{i,c}(n)\le G_i(\lambda)\log_2(n+c),
 \qquad G_i(z)=\sum_{j=0}^{i-1}z^j.
\end{equation}
In fact the shifted peak grows by at most the factor \((3/2)^{a_i}\)
through run \(i\). Hence
\begin{equation}\label{eq:peak-budget}
 H_{J,c}(n)+c\le(n+c)(3/2)^{A_{J,c}(n)} .
\end{equation}

For a fixed \(J\ge2\), abbreviate \(A_i=A_{i,c}(n)\) and write
\(V_i=v_1+\dots+v_i\) with \(V_0=A_0=0\). Set
\begin{align*}
 P&=3^{A_{J-1}}u_1, &
 Q&=2^{A_J-a_1+V_{J-1}}u_J,\\
 T_i&=c\,3^{A_{J-1}-A_i}2^{A_i-a_1+V_{i-1}}(2^{v_i}-1)
       && (1\le i<J).
\end{align*}
Multiplying \eqref{eq:transition} by
\(3^{A_{J-1}-A_i}2^{A_i-a_1+V_{i-1}}\) and telescoping gives
\begin{equation}\label{eq:telescope}
 Q=P+\sum_{i=1}^{J-1}T_i,\qquad
 P+\sum_{h<i}T_h
   =3^{A_{J-1}-A_{i-1}}2^{A_i-a_1+V_{i-1}}u_i .
\end{equation}
Since \(3^{a_i}u_i=2^{v_i}x_{i+1}+c\), the summands satisfy
\begin{equation}\label{eq:ratio}
 0<\frac{T_i}{P+\sum_{h<i}T_h}
   =\frac{c(2^{v_i}-1)}{3^{a_i}u_i}
   <\frac{c}{x_{i+1}}.
\end{equation}
In particular \(Q\le (1+c)^{J-1}P\).

Normalize the \(p\)-adic absolute values by \(|p|_p=p^{-1}\)
and use the ordinary absolute value at infinity. The local product is
\begin{equation}\label{eq:product}
 P|P|_3|Q|_2
 \prod_{i=1}^{J-1}\bigl(T_i|T_i|_2|T_i|_3\bigr)
 \le c^{J-1}(n+c)2^{-A_J}.
\end{equation}
Indeed \(P|P|_3\le u_1\), \(|Q|_2=2^{-A_J+a_1-V_{J-1}}\), and
each factor \(T_i|T_i|_2|T_i|_3\) is the part of \(c(2^{v_i}-1)\)
coprime to \(6\). It is a positive integer at most \(c2^{v_i}\).
The factors \(2^{v_i}\) cancel \(2^{-V_{J-1}}\), leaving
\(2^{a_1}u_1=n+c\).

\subsection{Projection on a minimal subspace}

We use the \(S\)-adic Subspace Theorem over \(\Q\)
\cite{Schlickewei}\cite[Theorem~2.3]{Bilu}.
In the integral-vector formulation, take
\(S=\{\infty,2,3\}\), and at each \(v\in S\) fix \(d\) independent
linear forms in \(d\) variables with rational coefficients.
For any \(\delta>0\), the nonzero \(y\in\Z^d\) with
\[
 \prod_{v\in S}\prod_{j=1}^{d}|L_{v,j}(y)|_v
       <\|y\|_\infty^{-\delta}
\]
lie in finitely many proper rational linear subspaces. Here
\(\|y\|_\infty\) is the maximum of the ordinary absolute coordinates.

\begin{lemma}\label{lem:projection}
Fix \(\ell\ge1\). Suppose \(P,T_1,\ldots,T_\ell\) are positive
integers along a sequence, \(P\to\infty\), and \(T_i/P\to0\) for
each \(i\). Put \(Q=P+\sum_iT_i\). For every \(\delta>0\) there
are only finitely many terms of the sequence satisfying
\[
 P|P|_3|Q|_2\prod_{i=1}^{\ell}\bigl(T_i|T_i|_2|T_i|_3\bigr)
     < P^{-\delta}.
\]
\end{lemma}

\begin{proof}
Suppose there are infinitely many such terms. Choose a rational subspace
\(W\) of least dimension containing infinitely many of them, and pass
to those terms. Dividing the vectors by \(P\) and taking the limit
shows that \(W\) contains \(e_P=(1,0,\ldots,0)\). Thus
\(W=\Q e_P\oplus W_T\), where \(W_T\) is the subspace of \(W\)
with \(P\)-coordinate zero. Choose indices \(I\) for which the
coordinate projection onto \((T_i)_{i\in I}\) is an isomorphism on
\(W_T\). Since all \(T_i>0\) on the sequence, \(I\) is nonempty.

On \(W\), work in the integral coordinates \(y=(P,(T_i)_{i\in I})\)
of the projected points. At infinity and at \(3\), take the coordinate
forms. At \(2\), take \(Q,(T_i)_{i\in I}\); they are independent since
the coefficient of \(P\) in \(Q\) on \(W\) is one. The new local
product is obtained by dropping the full factors
\(T_i|T_i|_2|T_i|_3\) for \(i\notin I\), each a positive integer.
Thus it remains less than \(P^{-\delta}\). For all sufficiently large
terms, \(\|y\|_\infty=P\). The Subspace Theorem puts the projections in
finitely many proper rational subspaces. Pulling these back under
the isomorphism gives a proper rational subspace of \(W\) containing
infinitely many of the original vectors, a contradiction.
\end{proof}

\begin{theorem}[Fixed-run odd-step budget]\label{thm:budget}
For every positive odd integer \(c\), integer \(J\ge1\), and real
\(\eta>0\), there is a constant \(C_{c,J,\eta}\) such that for every
positive odd \(n\),
\[
 A_{J,c}(n)\le(1+\eta)\log_2(n+c)+C_{c,J,\eta}.
\]
\end{theorem}

\begin{proof}
The case \(J=1\) follows from \(a_1=v_2(n+c)\). Induct on \(J\);
if the bound fails, there is a sequence \(n\to\infty\) with
\(A_J>(1+\eta)\log_2(n+c)\). For each \(2\le i\le J\), apply
the induction hypothesis with \(\eta/2\) to the first \(i-1\)
runs. The remaining odd-step count \(A_J-A_{i-1}\) then tends to
infinity. By \eqref{eq:elementary}, applied from \(x_i\) with
\(J-i+1\) runs, it is at most
\(G_{J-i+1}(\lambda)\log_2(x_i+c)\). It follows that every
\(x_i\) for \(2\le i\le J\) tends to infinity.
By \eqref{eq:ratio} and \(Q\le(1+c)^{J-1}P\), this gives
\(T_i/P\to0\). Also \(P\ge2^{a_1}u_1=n+c\), and
\eqref{eq:elementary} gives \(P\le(n+c)^{1+\lambda G_{J-1}(\lambda)}\).
By \eqref{eq:product} the local product is at most
\(c^{J-1}(n+c)^{-\eta}\). With
\(\delta=\eta/[2(1+\lambda G_{J-1}(\lambda))]\), it is less than
\(P^{-\delta}\) for all sufficiently large terms. Lemma~\ref{lem:projection}
gives the required contradiction.
\end{proof}

The argument fixes both the dimension and the exponent \(\eta\)
before applying the Subspace Theorem. No effective value of
\(C_{c,J,\eta}\) is extracted.

\begin{corollary}[Optimal fixed-run power]\label{cor:height}
For fixed \(c,J\) and every \(\varepsilon>0\), all positive odd \(n\)
satisfy \(H_{J,c}(n)+c\ll_{c,J,\varepsilon}(n+c)^{\lambda+\varepsilon}\).
Moreover
\[
 \lim_{R\to\infty}
 \frac{\ln\bigl(\max_{1\le n\le R,\ n\ {\rm odd}}(H_{J,c}(n)+c)\bigr)}
      {\ln R}=\lambda.
\]
\end{corollary}

\begin{proof}
Use Theorem~\ref{thm:budget} in \eqref{eq:peak-budget} with
\(\eta=\varepsilon/(\lambda-1)\). For the lower bound take
\(n=2^k-c\) with \(k=\lfloor\log_2(R+c)\rfloor\), so that \(n\) is
positive odd for sufficiently large \(R\). Its first shifted peak
is \(3^k\ge((R+c)/2)^\lambda\).
\end{proof}

\begin{remark}[A logarithmic excess over the pure power]\label{rem:excess}
For \(c=1\) let \(a\ge2\), \(b=2\cdot3^{a-1}-1\),
\(u=(2^{b+1}-1)/3^a\), and \(n=2^a u-1\).
The identity \(v_3(4^j-1)=1+v_3(j)\), by lifting the exponent,
shows that \(u\) is an odd integer. The first peak is \(2^{b+1}-2\);
one even step reaches \(2^b-1\), giving \(H_{2,1}(n)+1=3^b\).
Since \((n+1)(3/2)^a=2^{b+1}-1\) and \(b\asymp\ln(n+1)\),
\[
 H_{2,1}(n)+1\asymp (n+1)^\lambda(\ln(n+1))^{\lambda-1}.
\]
Thus the exponent limit in Corollary~\ref{cor:height} does not
give a bound with \(\varepsilon=0\) for \(J\ge2\).
\end{remark}

\section{Long trajectories and cycles}\label{sec:asymptotic}

Write \(\tau_{J,c}(n)\) for the number of shortcut steps from \(n\)
through the even steps after its \(J\)-th odd run, ending at the
start of run \(J+1\). We can let \(J\) vary by iterating a height
bound for a single fixed block size.

\begin{corollary}[Block bound]\label{cor:blocks}
For every positive odd \(c\) and fixed \(h\ge1\) there exists
a constant \(D_{c,h}\) such that for every odd \(n\ge1\) and \(J\ge1\),
\[
 \tau_{J,c}(n)+\log_2(H_{J,c}(n)+c)
 \le D_{c,h}\,2^{J/h}\bigl(1+\log_2(n+c)\bigr).
\]
\end{corollary}

\begin{proof}
Partition the runs into blocks of length \(h\) and a possible shorter
final block. If \(L_j=\log_2(x_{1+jh}+c)\) is the starting logarithm
of block \(j\), Corollary~\ref{cor:height} with
\(\lambda+\varepsilon=2\) bounds the logarithm of every shifted state
in that block by \(2L_j+C_{c,h}\), for a nonnegative constant.
The recurrence \(L_{j+1}\le2L_j+C_{c,h}\) implies
\(L_j\le2^j(L_0+C_{c,h})\).

For any run within the block, both its start and its even peak lie
below the block's height bound. If its start is \(x\) and its peak is
\(y\), it takes \(a\le\log_2(x+c)\) odd steps followed by
\(v_2(y)\le\log_2(y+c)\) even steps. Thus the block has at most
\(2h(2L_j+C_{c,h})\) steps. Sum over \(\lceil J/h\rceil\) blocks,
and bound the largest shifted logarithm by the largest block bound.
\end{proof}

In particular, for each \(\theta>0\), choose \(h\) so large that
\((\ln 2)/h<\theta\). The elapsed time and the logarithm of the
height after \(J\) runs are then both
\(O_{c,\theta}(e^{\theta J}(1+\log_2(n+c)))\). This estimate is
uniform over the starting integer.

For a positive cycle of \(S\), let \(t\) be its least period, \(s\)
its number of odd inputs, \(M\) its number of maximal odd runs, and
\(m\) its smallest state. The \(M\) run starts are exactly the local
minima. Start at \(m\), which must be odd. The product of successive
ratios around the cycle gives
\begin{equation}\label{eq:cycle-gap}
 \Lambda=t\ln2-s\ln3=\sum_{x\ {\rm odd\ in\ cycle}}
             \ln\!\left(1+\frac1{3x}\right)>0.
\end{equation}
Within a run from \(x_i\), the \(j\)-th odd input is at least
\((3/2)^j x_i\). Thus the run contributes strictly less than
\(\frac1{3x_i}\sum_{j\ge0}(2/3)^j=1/x_i\). This yields
\begin{equation}\label{eq:sum-min}
 \Lambda<\sum_{i=1}^M\frac1{x_i}\le M/m .
\end{equation}
In the form of Simons and de Weger \cite[Lemma~12]{SimonsWeger},
Rhin's inequality for a nontrivial positive cycle reads
\begin{equation}\label{eq:rhin}
 \Lambda>\exp\bigl(-13.3(0.46057+\ln s)\bigr).
\end{equation}
In particular \(m<2^9 M s^{13.3}\) and, since \(M\le s\),
\(\log_2(m+1)=O(1+\ln s)\).

\begin{corollary}\label{cor:superlog}
For every \(D>0\) there exists \(s_D\) such that every nontrivial
positive cycle with \(s\ge s_D\) satisfies \(M\ge D\ln s\).
\end{corollary}

\begin{proof}
Start at the minimum \(m\), so that \(\tau_{M,1}(m)=t\ge s\).
For every fixed \(h\ge1\), Corollary~\ref{cor:blocks} gives
\[
 s\ll_h 2^{M/h}(1+\log_2(m+1)).
\]
Use \eqref{eq:rhin} and take logarithms to get
\(
 M\ge (h/\ln 2)\bigl(\ln s-\ln\ln s-O_h(1)\bigr)
\)
for \(s\ge 2\). Choosing \(h>D\ln 2\) before letting \(s\) grow
proves the claim.
\end{proof}

\section{An exact certificate for \texorpdfstring{\(M\ge108\)}{M >= 108}}
\label{sec:certificate}

\begin{theorem}\label{thm:108}
Every nontrivial positive cycle of \(S\) has at least \(108\)
local minima.
\end{theorem}

This section uses Ba\v{r}ina's verification for \(n<2^{71}\)
\cite{Barina} and \eqref{eq:rhin}. Write \(A=2^{71}\);
since \(A\) itself is a power of two, a nontrivial cycle has \(m>A\).
Assume \(M\le107\). The argument covers all \(s\ge1\) by a tail
inequality, six unimodular brackets, and bounds on successive runs.

\subsection{The tail and the unimodular brackets}

Put \(U=1584962501/10^9>\lambda\), \(\mu=133/10\), and
\(Q_0=5225419486637257521420827\). By \eqref{eq:elementary}, \eqref{eq:sum-min},
and \eqref{eq:rhin},
\begin{equation}\label{eq:tail}
 s<G_M(U)\bigl(\log_2 M+10+\mu\log_2 s\bigr).
\end{equation}
The verifier checks the reverse inequality at \(Q_0\) with
\(M=107\), using binary lengths as upper bounds for logarithms.
The function \((\log_2 107+10+\mu\log_2 s)/s\) is decreasing for
\(s\ge Q_0\) by differentiation. Hence \(s<Q_0\).

\begin{lemma}\label{lem:bracket}
Let \(u=(p,q)\), \(v=(r,w)\) have positive integer entries, with
\(pw-qr=1\) and \(\Lambda(u)>0>\Lambda(v)\), where
\(\Lambda(t,s)=t\ln2-s\ln3\).
Every integer pair \((t,s)\) with \(0<s<q+w\) and \(\Lambda(t,s)>0\)
satisfies \(\Lambda(t,s)\ge\Lambda(u)\).
\end{lemma}

\begin{proof}
Write \((t,s)=\alpha u+\beta v\) with \(\alpha,\beta\in\Z\).
If \(\alpha\le0\), positivity of \(\Lambda(t,s)\) forces
\(\beta<0\), contradicting \(s>0\). Thus \(\alpha\ge1\).
Then \(\beta\ge1\) would give \(s\ge q+w\).
Therefore \(\beta\le0\), proving the bound.
\end{proof}

The six pairs in Table~\ref{tab:brackets} give the indicated
strict lower bounds \(g\) on \(A\Lambda(t,s)\) when \(0<s<q+w\).
The verifier checks determinants and signs by rational enclosures
of \(\ln2\) and \(\ln3\). For \(X=2,3\) and \(z=(X-1)/(X+1)\)
it sums \(2z^{2j+1}/(2j+1)\) for \(0\le j<120\), with positive
remainder at most \(2z^{241}/(241(1-z^2))\).

\begin{table}[htbp]
\centering
\small
\begin{tabular}{@{}c r r l@{}}
\toprule
Bracket & \multicolumn{1}{c}{\(u=(p,q)\)} &
 \multicolumn{1}{c}{\(v=(r,w)\)} & \(g\)\\
\midrule
1 & \(630118245525664765\) & \(6724555128221608268\) & \(358\)\\
  & \(397560349370386783\) & \(4242721909926539673\) &\\[2pt]
2 & \(7354673373747273033\) & \(36143248623210700400\) & \(63\)\\
  & \(4640282259296926456\) & \(22803850947114245497\) &\\[2pt]
3 & \(43497921996957973433\) & \(79641170620168673833\) & \(239/10\)\\
  & \(27444133206411171953\) & \(50247984153525417450\) &\\[2pt]
4 & \(123139092617126647266\) & \(202780263237295321099\) & \(803/100\)\\
  & \(77692117359936589403\) & \(127940101513462006853\) &\\[2pt]
5 & \(49373105075258054570781\) & \(2727782575569043909543559\) &
  \(9/10000\)\\
  & \(31150961018190238869556\) & \(1721039188200292347893905\) &\\[2pt]
6 & \(2777155680644301964114340\) & \(5504938256213345873657899\) &
  \(46/100000\)\\
  & \(1752190149218482586763461\) & \(3473229337418774934657366\) &\\
\bottomrule
\end{tabular}
\caption{For each pair, the numerator coordinate is on the first row
and the denominator coordinate on the second.}
\label{tab:brackets}
\end{table}

Let \(q_{37},q_{39},q_{41},q_{43},q_{47},Q_0\) be the successive
values \(q+w\) in the six rows. Bracket 1 excludes \(s<q_{37}\),
since \(A\Lambda<M\le107\). For the two highest intervals,
brackets 5 and 6 give \(m<107A/g\) by \eqref{eq:sum-min}, hence
\(m<2^{88}\) and \(m<2^{89}\), respectively.

\subsection{Low minima and odd-step bounds}

For a threshold \(B=2^h A\), let \(L_B\) be the number of run starts
at most \(B\). By \eqref{eq:sum-min},
\begin{equation}\label{eq:low}
 g<A\Lambda<L_B+(M-L_B)2^{-h}.
\end{equation}
For the first three intervals in Table~\ref{tab:bands}, this implies
the listed bound \(L_B\ge b\). The verifier checks
\((b-1)+(107-b+1)2^{-h}\le g\) in each case.

\begin{lemma}\label{lem:split}
Suppose at least \(b\ge1\) run starts in a cycle are at most \(B\).
Let \(F_d\) bound the number of odd steps in \(d\) successive runs
from any start at most \(B\), with \(F_0=0\) and nonnegative,
nondecreasing increments \(F_d-F_{d-1}\). If \(M\le107\), then
\[
 s\le F_{108-b}+(b-1)F_1 .
\]
\end{lemma}

\begin{proof}
Choose \(b\) of the starts at most \(B\). They partition the cycle
into \(b\) arcs of positive run lengths \(d_i\) with sum \(M\).
Moving one run from an arc of length at least two to a longest arc
can only increase \(\sum_i F_{d_i}\), by the increment assumption.
Iterate until all but one length equal one. Then replace \(M\) by
\(107\) in the resulting bound.
\end{proof}

For the first two intervals, use \(F_d=(71+h)G_d(U)\) by
\eqref{eq:elementary}, since an odd start \(x\le B\) has \(x+1\le B\).
The next two intervals use the shared two-run budget in
Section~\ref{sec:caps}. For the final interval, use the elementary
budget from the global minimum \(m<2^{89}\).

\begin{table}[htbp]
\centering
\begin{tabular}{@{}l r c r ll@{}}
\toprule
Interval for \(s\) & \(g\) & \(B\) & \(b\) &
Budget from \(B\) & Tested upper bound\\
\midrule
\([q_{37},q_{39})\) & \(63\) & \(8A\) & 57 & elementary &
  \(F_{51}+56F_1<q_{37}\)\\
\([q_{39},q_{41})\) & \(239/10\) & \(128A\) & 24 & elementary &
  \(F_{84}+23F_1<q_{39}\)\\
\([q_{41},q_{43})\) & \(803/100\) & \(4096A\) & 9 & two-run &
  \(F_{99}+8F_1<q_{41}\)\\
\([q_{43},q_{47})\) & \(9/10000\) & \(2^{88}\) & 1 & two-run &
  \(F_{107}<q_{43}\)\\
\([q_{47},Q_0)\) & \(46/100000\) & \(2^{89}\) & 1 & elementary &
  \(F_{107}<q_{47}\)\\
\bottomrule
\end{tabular}
\caption{Each tested upper bound contradicts the left endpoint.}
\label{tab:bands}
\end{table}

\subsection{Two-run caps and the logarithmic margin}\label{sec:caps}

Consider odd \(x<2^k\) and write \(x+1=2^a u\) with \(u\)
positive odd. The first peak is \(y=3^a u-1\).
Let \(v=v_2(y)\) and \(z=y/2^v\); set \(b=v_2(z+1)\).
The second shifted peak is \(3^b(z+1)/2^b\).
For a target \(K>k\), the first shifted peak is smaller than
\(2^K\) if \(3^k<2^K\), which the verifier checks.
If the second shifted peak reaches \(2^K\), then
\[
 2^K\le(3/2)^b(z+1)\le 2^{k+1-v}(3/2)^{a+b}.
\]
The last inequality follows from \(z\ge1\) and \(y+1\le2^k(3/2)^a\).
Since \(\lambda<U\), any such \(b\) is at least
\[
 b_0=\left\lceil\frac{K-k-1+v}{U-1}\right\rceil-a .
\]
Since \(y<2^{\,k+(U-1)a}\), put \(\ell=k+\lceil(U-1)a\rceil\).
We have \(1\le v<\ell\) and \(b\le\ell-v\). The condition
\(b_0\le\ell-v\) yields the integer upper cutoff
\[
 v\le
 \left\lfloor\frac{(U-1)(a+\ell)-(K-k-1)}{U}\right\rfloor .
\]
Since \(\ell\le k+(U-1)a+1\), feasibility with \(v\ge1\)
also requires \(a\ge (K-Uk)/(U(U-1))\). This allows the
verifier to start the search at a positive integer lower bound
on \(a\) when the target has a large margin.
Put \(d=\max(1,k-a)\). Because \(u\) is odd, the size condition
\(u\le2^{k-a}\) is equivalent to \(0<u<2^d\). The verifier
checks \(b_0>0\) and \(w=v+b_0\ge d\) at \(v=1\).
Both inequalities persist for larger \(v\).
By \eqref{eq:transition}, a qualifying \(u\) satisfies
\begin{equation}\label{eq:tworun-congruence}
 u\equiv I-I2^v\pmod{2^w},\qquad I\equiv 3^{-a}\pmod{2^w}.
\end{equation}
Conversely this congruence implies \(v_2(3^a u-1)=v\)
and \(v_2(z+1)\ge b_0\). Since \(w\ge d\), there is at most
one qualifying \(u\) per \(a,v\). It is the positive odd
representative in \([0,2^w)\) if it lies below \(2^d\).

\paragraph{A binary window test.}
Use a common inverse \(I=3^{-a}\bmod 2^L\) with \(L\) an odd integer
at least \(\max(\lceil Uk\rceil,k+32)\). This suffices since
\(w\le\ell\le\lceil Uk\rceil\). The lowest \(v\) bits of
\(I-I2^v\) equal those of \(I\). Thus, if the first set bit
of \(I\) at or above position \(d\) is at position \(j\),
we must have \(v\le j\). Positions are numbered from zero.

When \(v\le d\) and \(w-d\ge32\), the 32 bits of \(u\) starting
at position \(d\) must vanish. Let \(T\) be the 32 bits of \(I\)
starting at \(d\), read as an integer, and \(R\) the 32 bits
starting at \(d-v\). Binary subtraction gives
\[
 T-R-\beta\equiv0\pmod{2^{32}},\qquad \beta\in\{0,1\},
\]
where \(\beta\) is the borrow from lower positions. Therefore
\(R\) must equal \(T\) or \(T-1\) modulo \(2^{32}\). To screen
many \(v\) together, the verifier looks for the first three
whole bytes of these candidate windows, allowing all eight
bit alignments. More explicitly, if the window start has bit
offset \(r\) within a byte, the three bytes begin
\((-r)\bmod8\) positions inside the window, so all 24 tested
bits lie among the 32 required bits. Surviving candidates are
tested by the full congruence \eqref{eq:tworun-congruence}.
Small \(v\) with insufficient window width and larger \(v>d\)
are tested directly. All operations in this filter are exact.

\paragraph{The required margin.}
Writing \(k=\lceil\log_2(x+1)\rceil\) along the family of
Remark~\ref{rem:excess} gives
\[
 \log_2(H_{2,1}(x)+1)=\lambda k+(\lambda-1)\log_2 k+O(1).
\]
Thus a fixed additive allowance cannot cover all scales.
For the finite scale \(k=15416\), take
\(a=2884\), \(v=1\), and \(u=-3^{-a}\bmod 2^{12549}\)
with the least nonnegative representative. The appendix
checks \(x=2^a u-1<2^{15416}\), \(b=12548\), and
\(H_{2,1}(x)+1\ge2^{24441}\). Here
\(\lceil Uk\rceil=24434\), so the allowance \(5\) (and even
\(7\)) fails. The logarithmic term has size approximately
\(8.14\) at this scale. The full residue test certifies the
target \(24442\), so \(8\) is the least integer allowance for
this \(k\). The iteration below uses the certified allowance \(10\).

\paragraph{Iteration.}
Starting at \(k_0=88\), take the seventeen successive caps
in Table~\ref{tab:caps}. For \(0\le j<17\) the verifier confirms
\(H_{2,1}(x)+1<2^{k_{j+1}}\) for all positive odd \(x<2^{k_j}\).
The next start after the two runs is therefore below \(2^{k_{j+1}}\).
The lengths of the two runs are bounded by \(k_j\) and \(Uk_j\),
respectively, by \eqref{eq:elementary}. After these seventeen pairs,
bound individual run lengths by \(k_{17},Uk_{17},U^2 k_{17},\ldots\).
The combined list is nondecreasing because \(k_{j+1}>Uk_j\).
Taking \(F_d\) as the sum of its first \(d\) entries proves the
budget needed for the third and fourth rows of Table~\ref{tab:bands}.
The verifier confirms both inequalities. In the final cap the
available slack in \(F_{107}<q_{43}\) permits a larger target,
which raises the lower cutoff on \(a\) and speeds this check.
Table~\ref{tab:caps}
gives the cap chain and the number of residues retained after
the size restriction on \(u\). This exhausts all \(s\) and proves
Theorem~\ref{thm:108}.

\begin{table}[htbp]
\centering
\begin{tabular}{@{}r rr r@{}}
\toprule
\(j\) & \(k_j\) & \(k_{j+1}\) & Residues tested\\
\midrule
0 & 88 & 145 & 0\\
1 & 145 & 235 & 0\\
2 & 235 & 378 & 0\\
3 & 378 & 605 & 0\\
4 & 605 & 964 & 1\\
5 & 964 & 1533 & 1\\
6 & 1533 & 2435 & 0\\
7 & 2435 & 3865 & 5\\
8 & 3865 & 6131 & 9\\
9 & 6131 & 9723 & 0\\
10 & 9723 & 15416 & 5\\
11 & 15416 & 24444 & 0\\
12 & 24444 & 38753 & 0\\
13 & 38753 & 61428 & 6\\
14 & 61428 & 97372 & 0\\
15 & 97372 & 154343 & 0\\
16 & 154343 & 301000 & 0\\
\bottomrule
\end{tabular}
\caption{For each row, odd seeds below \(2^{k_j}\) have shifted
two-run peak less than \(2^{k_{j+1}}\).}
\label{tab:caps}
\end{table}

\section{Scope of the bounds}

The odd-step budget controls every prescribed finite number of runs.
Its constants are obtained after fixing the number of runs, so the
estimate does not control an entire infinite trajectory. The
superlogarithmic consequence follows by fixing a block length first.
The finite cycle bound uses a separate effective argument, and its
arithmetic checks do not invoke any constants from the Subspace Theorem.

\appendix
\section{Complete arithmetic verifier}\label{app:code}

Save the listing as \texttt{certificate.py} and run it with Python
3.8 or later, with assertions enabled. The standard-library verifier
uses up to twelve worker processes by default. Pass \texttt{1} as the
first argument to use one process. The final line confirms that the
six brackets, the tail, the logarithmic-margin example, and all
seventeen two-run caps passed. The external inputs are
Ba\v{r}ina's verification threshold and the explicit Rhin bound
in \eqref{eq:rhin}.

\lstset{language=Python,basicstyle=\small\ttfamily,
  columns=fullflexible,keepspaces=true,showstringspaces=false,
  breaklines=true,aboveskip=0.7\baselineskip}
\begin{lstlisting}
from concurrent.futures import ProcessPoolExecutor
from fractions import Fraction as F
from itertools import chain
from os import cpu_count
from sys import argv
from time import perf_counter


A = 2**71
M0 = 107
P, D = 1584962501, 10**9
N = P - D
U = F(P, D)
MU = F(133, 10)
WIDTH = 32

CAPS = (88, 145, 235, 378, 605, 964, 1533, 2435, 3865, 6131, 9723,
        15416, 24444, 38753, 61428, 97372, 154343, 301000)

BRACKETS = (
    ((630118245525664765, 397560349370386783),
     (6724555128221608268, 4242721909926539673), F(358)),
    ((7354673373747273033, 4640282259296926456),
     (36143248623210700400, 22803850947114245497), F(63)),
    ((43497921996957973433, 27444133206411171953),
     (79641170620168673833, 50247984153525417450), F(239, 10)),
    ((123139092617126647266, 77692117359936589403),
     (202780263237295321099, 127940101513462006853), F(803, 100)),
    ((49373105075258054570781, 31150961018190238869556),
     (2727782575569043909543559, 1721039188200292347893905), F(9, 10000)),
    ((2777155680644301964114340, 1752190149218482586763461),
     (5504938256213345873657899, 3473229337418774934657366), F(46, 100000)),
)


def ceil_div(x, y):
    return (x + y - 1)//y


def log_box(x, terms=120):
    z = F(x-1, x+1)
    lo = sum(2*z**(2*j+1)/(2*j+1) for j in range(terms))
    tail = 2*z**(2*terms+1)/((2*terms+1)*(1-z*z))
    return lo, lo+tail


def geom(j):
    return (U**j-1)/(U-1)


def odd_budget(j):
    entries = list(chain.from_iterable((k,U*k) for k in CAPS[:-1]))
    assert all(x <= y for x,y in zip(entries, entries[1:]+[CAPS[-1]]))
    excess = max(0,j-len(entries))
    return sum(entries[:j]) + CAPS[-1]*geom(excess)


def cycle_arithmetic():
    lo2, hi2 = log_box(2)
    lo3, hi3 = log_box(3)
    assert hi3/lo2 < U and MU*F(46057,100000) < 9*lo2
    stops = []
    for (p,q),(r,w),g in BRACKETS:
        assert p*w-q*r == 1
        assert A*(p*lo2-q*hi3) > g and A*(r*hi2-w*lo3) < 0
        stops.append(q+w)
    q37,q39,q41,q43,q47,end = stops
    assert stops == sorted(stops) and BRACKETS[0][2] > M0
    assert end > geom(M0)*(M0.bit_length()+10+MU*end.bit_length())

    for i,(h,b,left) in enumerate(
        ((3,57,q37),(7,24,q39),(12,9,q41)), start=1
    ):
        length = M0-b+1
        assert b-1+F(length,2**h) <= BRACKETS[i][2]
        if i < 3:
            assert (71+h)*(geom(length)+b-1) < left
        else:
            assert 71+h < CAPS[0]
            assert odd_budget(length)+(b-1)*CAPS[0] < left
    assert M0 < BRACKETS[4][2]*2**(88-71)
    assert odd_budget(M0) < q43
    assert M0 < BRACKETS[5][2]*2**(89-71)
    assert 89*geom(M0) < q47


def window_candidates(inverse, d, low, high):
    span = high-low+WIDTH
    band = (inverse >> (d-high)) & ((1 << span)-1)
    word = band.to_bytes(ceil_div(span,8),"little")
    target = (inverse >> d) & ((1 << WIDTH)-1)
    for residue in range(min(8,high-low+1)):
        shift = (-residue)%8
        start = (residue+shift)//8
        end = start+(high-low-residue)//8+3
        patterns = {((target-borrow) >> shift) & 0xffffff
                    for borrow in (0,1)}
        for value in patterns:
            pattern = value.to_bytes(3,"little")
            at = word.find(pattern,start,end)
            while at >= 0:
                yield high+shift-8*at
                at = word.find(pattern,at+1,end)


def check_interval(job):
    k, target, begin, end = job
    bits = max(ceil_div(P*k,D),k+WIDTH) | 1
    modulus = 1 << bits
    one = (modulus+1)//3
    offsets = (0,one,2*one)
    inverse = pow(3,1-begin,modulus)
    need = [ceil_div(D*(target-k-1+v),N) for v in range(bits+1)]
    full = 1
    while need[full]+full-k-1 < WIDTH:
        full += 1
    checked = 0
    for a in range(begin,end):
        inverse, remainder = divmod(inverse,3)
        inverse += offsets[remainder]
        d = max(1,k-a)
        l = k+ceil_div(N*a,D)
        hi = min(l-1,(N*(a+l)-D*(target-k-1))//P)
        if hi < 1:
            continue
        assert need[1]-a > 0 and need[1]-a+1 >= d
        high_bits = inverse >> d
        if high_bits:
            first = (high_bits & -high_bits).bit_length()-1
            hi = min(hi,d+first)
        middle = min(d,hi)
        windows = (window_candidates(inverse,d,full,middle)
                   if full <= middle else ())
        small = range(1,min(full-1,hi)+1)
        large = range(max(full,d+1),hi+1)
        for v in chain(small,windows,large):
            w = v+need[v]-a
            u = (inverse-(inverse << v)) & ((1 << w)-1)
            if u.bit_length() > d:
                continue
            checked += 1
            z1 = ((3**a*u-1) >> v)+1
            assert z1 % 2 == 0
            b = (z1 & -z1).bit_length()-1
            assert 3**b*(z1 >> b) < 1 << target, (k,target,a,v,b)
    return checked


def witness():
    a,v,b = 2884,1,12548
    mod = 1 << (v+b)
    u = (-pow(3,-a,mod)) % mod
    z1 = ((3**a*u-1) >> v)+1
    assert z1 % (1 << (b+1)) == 1 << b
    assert 0 < ((u << a)-1) < 1 << 15416
    assert 3**b*(z1 >> b) >= 1 << 24441
    assert 3**15416 < 1 << 24442
    assert check_interval((15416,24442,1,15417)) == 2


def check_caps(pool, workers):
    for k,target in zip(CAPS,CAPS[1:]):
        assert 3**k < 1 << target and target > U*k
        start = max(1,ceil_div((target*D-P*k)*D,P*N))
        if pool is None or k < 2000:
            count = check_interval((k,target,start,k+1))
        else:
            chunk = max(1,ceil_div(k-start+1,3*workers))
            jobs = ((k,target,a,min(a+chunk,k+1))
                    for a in range(start,k+1,chunk))
            count = sum(pool.map(check_interval,jobs))
        print("Two-run cap",k,target,"residues tested",count,flush=True)


if __name__ == "__main__":
    start = perf_counter()
    cycle_arithmetic()
    witness()
    workers = int(argv[1]) if len(argv) > 1 else min(12,cpu_count() or 1)
    assert workers > 0
    if workers == 1:
        check_caps(None,1)
    else:
        with ProcessPoolExecutor(max_workers=workers) as pool:
            check_caps(pool,workers)
    print(f"All s covered for M <= {M0} in {perf_counter()-start:.2f} s.")
\end{lstlisting}

\begin{thebibliography}{99}

\bibitem{Alick}
J. Alick, \emph{Collatz cycles with at most 101 local minima:
a Lean 4 formalization}, public code repository,
\url{https://github.com/justinalick/collatz-m-cycles-101-lean}.

\bibitem{Barina}
D. Ba\v{r}ina, \emph{Improved verification limit for the
convergence of the Collatz conjecture},
J. Supercomput. \textbf{81} (2025), article 810.
\href{https://doi.org/10.1007/s11227-025-07337-0}{doi:10.1007/s11227-025-07337-0}.

\bibitem{Bilu}
Y. Bilu, \emph{The many faces of the Subspace Theorem
(after Adamczewski, Bugeaud, Corvaja, Zannier\ldots)},
Ast\'erisque \textbf{317} (2008), Exp. No. 967, 1--38.
\url{https://www.numdam.org/item/AST_2008__317__1_0/}.

\bibitem{Brox}
T. Brox, \emph{Collatz cycles with few descents},
Acta Arith. \textbf{92} (2000), 181--188.
\href{https://doi.org/10.4064/aa-92-2-181-188}{doi:10.4064/aa-92-2-181-188}.

\bibitem{Hercher}
C. Hercher, \emph{There are no Collatz \(m\)-cycles with \(m\le91\)},
J. Integer Seq. \textbf{26} (2023), Article 23.3.5.
\url{https://cs.uwaterloo.ca/journals/JIS/VOL26/Hercher/hercher5.html}.

\bibitem{Schlickewei}
H. P. Schlickewei, \emph{On linear forms with algebraic coefficients
and diophantine equations}, J. Number Theory \textbf{9} (1977),
381--392.
\href{https://doi.org/10.1016/0022-314X(77)90073-7}{doi:10.1016/0022-314X(77)90073-7}.

\bibitem{SimonsWeger}
J. Simons and B. de Weger, \emph{Theoretical and computational bounds
for \(m\)-cycles of the \(3n+1\) problem}, Acta Arith. \textbf{117}
(2005), 51--70.
\href{https://doi.org/10.4064/aa117-1-3}{doi:10.4064/aa117-1-3}.

\bibitem{Weaver}
L. Weaver, \emph{Fixed-descent excursion bounds and
superlogarithmic non-one counts}, public research note, 21 October
2026, \url{https://rstation.ai/r/collatz/posts/23}.

\bibitem{Wang}
X. Wang, \emph{Non-existence of Collatz \(m\)-cycles for \(m\le95\)},
preprint, 2026. \href{https://doi.org/10.5281/zenodo.21017296}{Zenodo 21017296}.

\end{thebibliography}
\end{document}
